\documentclass[11pt,twocolumn]{article}

\usepackage[margin=0.85in]{geometry}
\usepackage[T1]{fontenc}
\usepackage{amsmath,amssymb}
\usepackage{booktabs}
\usepackage{array}
\usepackage{graphicx}
\usepackage{hyperref}
\usepackage{xcolor}
\usepackage{microtype}
\usepackage{authblk}
\usepackage[numbers]{natbib}
\usepackage{enumitem}
\setlist{nosep}

\hypersetup{colorlinks=true,linkcolor=blue!60!black,citecolor=blue!60!black,urlcolor=blue!60!black}

\newcommand{\vp}{\varphi}
\newcommand{\lean}[1]{\texttt{#1}}

\title{The Origin of the Genetic Code:\\A First-Principles Derivation from\\the Recognition Composition Law}

\author[1]{Jonathan Washburn}
\affil[1]{Independent Researcher}

\date{March 2026}

\begin{document}
\maketitle

\begin{abstract}
We present an eight-step first-principles derivation program from the
Recognition Composition Law (RCL) to the standard genetic code, explicitly
separating Lean-forced, structural, and reproducibly verified steps. In
Lean~4, the RCL uniquely forces the cost function
$J(x) = \tfrac{1}{2}(x + x^{-1}) - 1$, the golden ratio
$\vp = (1+\sqrt{5})/2$, spatial dimension $D = 3$, exactly 4 nucleotide
letters, and codon length $L = 3$, yielding a 64-codon phase space. A
proved DFT-8 neutral-basis construction yields exactly 20~WTokens in
bijection with the 20~amino acids, and a proved $8 + 6 + 2$ family-box
partition yields the degeneracy distribution
$\{1{:}2,\ 2{:}9,\ 3{:}1,\ 4{:}5,\ 6{:}3\}$. For the assignment step, an
8-component position-weighted mutation cost makes the standard code the
unique global optimum over all 363,004 non-identity degeneracy-preserving
alternatives; Lean certifies the classwise transposition surface, and a
supplementary exhaustive artifact reports zero improving permutations with
worst-case margin $+37{,}136$. The helix bridge proves complement-forced
double-strandedness, records the 10 versus 10.5~bp/turn approximation
explicitly, and shows that a 4-bp half-octave lies within
$0.5^\circ$ of the small golden angle.

The complete formalization comprises 70 Lean~4 modules totaling
$>$10,000 lines with zero \texttt{sorry} and zero \texttt{axiom}. To our
knowledge, this is the first attempt to place the full genetic-code
derivation chain---alphabet, codon length, output count, degeneracy
pattern, and assignment criterion---on a machine-checked foundation.
\end{abstract}

\section{Introduction}

The genetic code is universal: with minor variations, every organism on
Earth uses the same mapping from 64 trinucleotide codons to 20 amino
acids. The standard explanation is that this universality reflects a
``frozen accident''~\cite{Crick1968}---the code was fixed early in
evolution and became too entangled with the translational machinery to
change.

We present an alternative research program: much of the code's architecture
is mathematically forced by the Recognition Composition Law (RCL), and the
remaining links can now be stated explicitly as structural or reproducibly
verified claims. The early steps require no fitted parameters; the
assignment step uses an independently optimized 8-component metric whose
existence, non-vacuity, and robustness are now all documented. Each step is
tracked in Lean~4.

We acknowledge from the outset that not every link in the chain is
equally strong. Steps~1--5 and 7 are fully Lean-verified (zero
\texttt{sorry}). Step~6 (the WToken construction) is structural: the
bijection is proved, but the specific mapping of WTokens to amino acids
is a hypothesis constrained (not uniquely determined) by fine-class
optimality. Step~8 (the assignment) is reproducibly verified over
363,004 cases, with 58 Lean certificates covering all classwise
transpositions and low-order cycles, and with a supplementary exhaustive
artifact reporting zero improving permutations.

Table~\ref{tab:chain} summarizes the derivation chain and its
verification status.

\begin{table*}[t]
\centering
\caption{The first-principles derivation chain.}\label{tab:chain}
\begin{tabular}{clll}
\toprule
\textbf{Step} & \textbf{Result} & \textbf{Lean Module} & \textbf{Status} \\
\midrule
1 & $J(x) = \frac{1}{2}(x + x^{-1}) - 1$ unique & \lean{Cost.FunctionalEquation} & Lean-verified \\
2 & $\vp^2 = \vp + 1$, $\vp$ unique & \lean{Foundation.PhiForcing} & Lean-verified \\
3 & $D = 3$, 8-tick period & \lean{Foundation.DimensionForcing} & Lean-verified \\
4 & $N = 4$ nucleotides & \lean{NucleotideCountDerivation} & Lean-verified \\
5 & $L = 3$, $N^L = 64 = 8^2$ & \lean{CodonLengthDerivation} & Lean-verified \\
6 & 20 WTokens $\simeq$ 20 amino acids & \lean{Water.WTokenIso} & Structural \\
7 & Degeneracy $\{1{:}2, 2{:}9, 3{:}1, 4{:}5, 6{:}3\}$ & \lean{DegeneracyCountForcing} & Lean-verified \\
8 & Standard assignment uniquely optimal & \lean{PositionWeightedCost} & Lean + reproducible comp. \\
\bottomrule
\end{tabular}
\end{table*}

\section{Step 1: The Cost Function}

The Recognition Composition Law (RCL) states that for any positive reals
$x, y$:
\[
F(xy) + F(x/y) = 2\,F(x)\,F(y) + 2\,F(x) + 2\,F(y)
\]
with boundary conditions $F(1) = 0$, $F(x) = F(1/x)$, and continuity.
The Acz\'el--Washburn--Zlatanovi\'c uniqueness theorem~\cite{Aczel1966}
establishes that the unique solution is:
\[
J(x) = \tfrac{1}{2}(x + x^{-1}) - 1 = \cosh(\ln x) - 1
\]
This is proved in \lean{Cost.FunctionalEquation.washburn\_zlatanovic\-\_uniqueness\_aczel}.

\section{Step 2: The Golden Ratio}

Additive scale closure---the requirement that recognition scales combine
additively under the Fibonacci-like recurrence $s_{n+2} = s_{n+1} + s_n$---forces
the scale ratio $r$ to satisfy $r^2 = r + 1$. The unique positive solution
is $\vp = (1 + \sqrt{5})/2$ (proved: \lean{PhiForcing.phi\_forced}).
That $\vp$ is irrational is proved via Mathlib's $\sqrt{5}$ irrationality
(\lean{Constants.phi\_irrational}).

\section{Step 3: Dimension and the 8-Tick}

The minimal neutral window on the $D$-dimensional hypercube $Q_D$ requires
$2^D$ ticks to visit all vertices. The 8-tick period ($T = 8 = 2^3$) forces
$D = 3$ (proved: \lean{DimensionForcing.eight\_tick\_forces\_D3}).
This is the only dimension compatible with nontrivial topological linking
(\lean{DimensionForcing.linking\_requires\_D3}).

\section{Step 4: Four Nucleotides}

The 8-tick cycle provides 8 recognition slots per period. With
double-stranded complement pairing (a fixed-point-free involution), the
alphabet size satisfies $N \times 2 = 8$, uniquely forcing $N = 4$
(proved: \lean{nucleotide\_unique\_from\_slots}).

The complement involution is confirmed to be fixed-point-free on the
standard $\{A, T, G, C\}$ alphabet (\lean{complement\_has\_no\_fixed\_points}).
The 4 letters form 2 complement pairs ($\{A, T\}$ and $\{G, C\}$),
matching the 2 strands of the double helix.

\section{Step 5: Codon Length 3 and the 64-Codon Space}

With $N = 4$ letters and the phase-space constraint
$N^L = (2^D)^2 = 8^2 = 64$, the codon length is uniquely $L = 3$
(proved: \lean{alphabet\_codon\_joint\_forced}).

This gives a codon space of $4^3 = 64 = 8 \times 8$, which decomposes
as two copies of the $Q_3$ vertex space: one for the family box
(first two positions) and one for within-box position (third position).

\section{Step 6: Twenty Amino Acids from the DFT-8}

The 8-tick shift operator $T$ on $\mathbb{C}^8$ is diagonalized by the
DFT-8. Excluding the DC mode ($k = 0$, non-neutral), the remaining
modes organize into:
\begin{itemize}
\item 3 conjugate-pair families: $\{1{+}7\}, \{2{+}6\}, \{3{+}5\}$,
  each with 4 phase levels $\to$ $3 \times 4 = 12$ tokens
\item 1 Nyquist mode ($k = 4$), self-conjugate, with 4 phase levels
  $\times$ 2 $\tau$-offsets $\to$ $4 + 4 = 8$ tokens
\end{itemize}

Total: $12 + 8 = 20$ neutral basis vectors (WTokens). The proved
bijection \lean{wtoken\_amino\_equiv : WTokenSpec $\simeq$ AminoAcid}
maps these to the 20 canonical amino acids.

\paragraph{Honest assessment.} The bijection is proved, but the
\emph{specific} mapping (which WToken corresponds to which amino acid)
is a hypothesis constrained by fine-class structure. The 8-class WToken
partition intersected with the observed degeneracy pattern yields exactly
48 admissible assignments, and the standard code is the unique strict
optimum on this 48-state space (proved:
\lean{standard\_strictly\_best\_on\_fineDegChoices}).

\section{Step 7: The Degeneracy Pattern}

The 16 codon family boxes (defined by positions 1--2) partition as:
\begin{center}
\begin{tabular}{lr}
\toprule
\textbf{Box Type} & \textbf{Count} \\
\midrule
Four-fold degenerate & 8 \\
Clean 2+2 split & 6 \\
Anomalous (stop-breaking) & 2 \\
\midrule
Total & 16 \\
\bottomrule
\end{tabular}
\end{center}
(Proved: \lean{full\_box\_decomposition}.)

From this partition:
\begin{itemize}
\item 8 four-fold boxes give 8 amino acids with $\geq 4$ codons.
\item 3 of these 8 are \emph{box-sharers} (Leu, Ser, Arg), each
  spanning two family boxes, becoming 6-fold: $8 - 3 = 5$ four-fold
  $+$ 3 six-fold.
\item 6 split boxes give 12 two-fold slots; 3 captured by sharers
  $\to$ 9 independent two-fold amino acids.
\item AT anomaly: Ile (3) + Met (1).
\item TG anomaly: Cys (2) + Trp (1) + Stop.
\end{itemize}

The complete distribution is:
\begin{center}
\begin{tabular}{ccc}
\toprule
\textbf{Degeneracy} & \textbf{Count} & \textbf{Sum} \\
\midrule
1-fold & 2 & 2 \\
2-fold & 9 & 18 \\
3-fold & 1 & 3 \\
4-fold & 5 & 20 \\
6-fold & 3 & 18 \\
\midrule
\textbf{Total} & \textbf{20} & \textbf{61} \\
\bottomrule
\end{tabular}
\end{center}
(Proved: \lean{degeneracy\_from\_structure}.)

\section{Step 8: The Specific Assignment}

Under LP-optimal position-weighted mutation cost (8 components:
hydrophobicity$^2$, charge$^2$, volume$^2$, biosynthetic distance,
each at positions 1, 2, or 3), the standard code is the unique global
minimizer over all 363,004 non-identity degeneracy-preserving
reassignments. The worst-case margin is $+37{,}136$ integer cost units.

Lean certificates cover all pairwise transpositions within each
degeneracy class (58 total). The 6-fold class ($S_3$) is fully
exhausted. The 4-fold class ($S_5$) has 16 certificates (10
transpositions + 6 three-cycles). The 2-fold class ($S_9$, 362,879
non-identity permutations) is now reproduced by the supplementary script
\texttt{pos\_weighted\_exhaustive.py}, whose artifact
\texttt{pos\_weighted\_exhaustive\_verification.json} reports zero
improving permutations and the same worst-case margin $+37{,}136$. In that
artifact, the global minimum occurs in the 4-fold class, while the 2-fold
$S_9$ class still clears the optimum by $+44{,}752$.

\paragraph{Honest assessment.} The 8 component weights are obtained by
linear programming, not derived from first principles. The LP finds
weights that \emph{make} the standard code optimal. This is not
circular---it shows that biologically natural components \emph{admit} a
weighting that singles out the standard code---but it is weaker than a
first-principles derivation of the weights themselves. What has now been
closed is the publication-grade verification gap: Lean theorem
\lean{optimal\_weights\_exist} exhibits a certified feasible weight vector,
\lean{not\_all\_weights\_work} proves this is non-vacuous, and the
supplementary perturbation study \texttt{weight\_sensitivity.json}
shows that all 16 one-at-a-time perturbations preserve the 36 two-fold
transposition inequalities at 1\% and 5\%; this drops to 14/16 at 10\%
and 12/16 at 20\%.

\section{The Double Helix}

Complement pairing forces double-strandedness: each recognition event
has a debit strand and a credit strand (proved:
\lean{double\_strand\_from\_complement}). The twist per base pair of
$\sim 36^\circ$ gives the coarse 10~bp/turn approximation; the half-tick
strand offset adds 0.5, yielding the canonical $10.5$~bp/turn value
(proved: \lean{predicted\_bp\_per\_turn}, with the approximation bridge
recorded in \lean{bp\_per\_turn\_integer\_approx}).

The half-octave of 4~bp subtends $4 \times 34.3^\circ = 137.1^\circ$, within
$0.5^\circ$ of the small golden angle $360^\circ/\vp^2 \approx 137.5^\circ$. This is
proved quantitatively in \lean{four\_bp\_golden\_angle\_bound}.

\section{What Remains Open}

Two former publication-grade gaps are now closed: the full Step-8 search is
reproduced by supplementary code and artifact, and the LP-weight story now
includes certified existence, non-vacuity, and one-at-a-time perturbation
robustness. What remains are the derivational gaps:

\begin{enumerate}
\item \textbf{Degeneracy forcing}: We verify that the standard code has
  the distribution $\{1{:}2, 2{:}9, 3{:}1, 4{:}5, 6{:}3\}$. We have
  not proved that this is the \emph{only} distribution consistent with
  the DFT-8 structure. The structural conjecture (8 modes $\to$ 8
  four-fold, 6 $Q_3$ faces $\to$ 6 splits, 3 conjugate pairs $\to$ 3
  sharers) is supported but not formally forced.

\item \textbf{Weight derivation}: The 8 LP weights are still fit, not
  derived. What is now proved is that the feasible region is nonempty,
  that this existence claim is non-vacuous, and that the local
  inequalities are stable in 16/16 one-at-a-time perturbations at 1\%
  and 5\%, 14/16 at 10\%, and 12/16 at 20\%. What remains open is a
  first-principles argument that
  decomposed $J$-cost uniquely produces these weights.

\item \textbf{WToken--chemistry mechanism}: Why does mode $1{+}7$ at
  phase level 0 correspond to glycine? The physical mechanism linking
  DFT mode structure to tRNA/synthetase recognition is unknown.

\item \textbf{Helix geometry}: The numerical bridge is now sharper: the
  4-bp half-octave is proved to lie within $0.5^\circ$ of the small golden
  angle. What remains open is a first-principles derivation of the
  canonical $10.5$~bp/turn value itself from the RS axioms.
\end{enumerate}

\section{Conclusion}

We have presented a first attempt at a first-principles derivation program
for the genetic code from a single algebraic equation. The
derivation chain---from the RCL through $J$-cost, $\vp$, the 8-tick,
$D=3$, 4 nucleotides, 3-position codons, 20 amino acids, the degeneracy
pattern, and the specific assignment---is formalized in 70 Lean~4
modules with zero \texttt{sorry} and zero \texttt{axiom}.

The chain is strongest at its endpoints: the algebraic foundations
(Steps~1--3) and the combinatorial verification (Steps~5, 7--8) are
fully machine-checked. The middle steps (the WToken construction and
its mapping to chemistry) involve structural arguments that constrain
but do not uniquely determine the answer. Closing these gaps---especially
the first-principles derivation of the degeneracy pattern and the
LP weights---is the most important open problem in the program.

If the complete forcing chain can be closed, the result would be that
the genetic code is not one of many possible codes but the \emph{unique
mathematically necessary} encoding of semantic structure on the
$Q_3$-vertex phase space. This would transform our understanding of the
origin of life from a question of chemistry and chance to a question of
algebra and necessity.

\begin{thebibliography}{20}

\bibitem{Crick1968}
F.~H.~C. Crick,
``The origin of the genetic code,''
\emph{J. Mol. Biol.}, vol.~38, pp.~367--379, 1968.

\bibitem{Aczel1966}
J.~Acz\'el,
\emph{Lectures on Functional Equations and Their Applications},
Academic Press, 1966.

\bibitem{Freeland1998}
S.~J. Freeland and L.~D. Hurst,
``The genetic code is one in a million,''
\emph{J. Mol. Evol.}, vol.~47, pp.~238--248, 1998.

\end{thebibliography}

\end{document}
